% SHARD-ID: CA-18-GALILEAN-SYMMETRY-MOTIVATION
% SHARD-TITLE: Why Galilean Symmetry Is a Different Lattice Target
% SHARD-SUMMARY: Opens the Galilean follow-up block by separating nonrelativistic spacetime symmetry from the Lorentzian energy-moment boost ansatz.
% SHARD-SUMMARY: Records the local source boundary: Galilei symmetry and projective-unitary lifting are sourced, while the named Bargmann mass extension remains a source gap.
% SHARD-SUMMARY: Identifies conserved mass or particle-number density as extra compiler input required for Galilean boost candidates.
% SHARD-KEYWORDS: Galilei group, Galilean symmetry, nonrelativistic symmetry, mass density, particle number, Bargmann source gap

\section{Why Galilean Symmetry Is a Different Lattice Target}
\label{sec:galilean-symmetry-motivation}

\subsection{Status, Sources, and Dependencies}

\textbf{Status: sourced boundary plus proposal shard.}  The Galilei group as a
classical nonrelativistic symmetry group is sourced locally.  The general
projective-representation and central-extension mechanism is sourced locally.
The specific mass central extension of the Galilei group is not yet sourced in
this repository; CA-20 treats its coefficient as a checked canonical-commutator
derivation, not as a literature claim about a named group.

Dependencies: CA-05 for projective symmetries via central extensions, CA-10--CA-17
for the Lorentzian lattice-symmetry programme, and CONVENTIONS.md (i).

% Source: references/cft/Schottenloher2008/Schottenloher2008.md:1440 -- examples of classical systems with a symmetry group.
% Source: references/cft/Schottenloher2008/Schottenloher2008.md:1443 -- Galilei group for free particles in classical nonrelativistic mechanics.
% Source: references/cft/Schottenloher2008/Schottenloher2008.md:1524 -- unitary representation of a topological group.
% Source: references/cft/Schottenloher2008/Schottenloher2008.md:1528 -- projective representation into U(P).
% Source: references/cft/Schottenloher2008/Schottenloher2008.md:1570 -- central extension by U(1) for a projective representation.
% Source: references/cft/Schottenloher2008/Schottenloher2008.md:1618 -- topological central-extension lift.
% Source: references/cft/Schottenloher2008/Schottenloher2008.md:1622 -- quantization as a central extension by U(1) of a covering group.

\subsection{The New Target}

The previous block targeted Lorentz-invariant continuum systems.  Its first
lattice intuition was:
\[
  \text{Lorentz boost} \simeq \text{position-dependent time translation},
  \qquad
  K^{\mathrm{Lor}}_a\sim\sum_x x_a h_x.
\]
That is appropriate only when the continuum target mixes space and time in the
Lorentzian way and the boost is controlled by the energy density.

The nonrelativistic target is different.  Schottenloher lists the Galilei group
as the symmetry group for free particles in classical nonrelativistic mechanics.
For the lab-book compiler this means that a Galilean target is not obtained by
renaming the Poincare block.  It has absolute time, spatial translations,
rotations, and boosts whose quantum implementation may be projective.

\subsection{Immediate Consequence for Lattices}

The Lorentzian boost ansatz used only the Hamiltonian density \(h_x\).  A
Galilean boost needs more input.  At time \(t=0\), the familiar nonrelativistic
one-particle expression is controlled by the mass-weighted position, not by the
energy-weighted position:
\[
  G_a(0)\sim mX_a.
\]
For a many-body lattice this suggests a mass-density or particle-number-density
first moment:
\[
  G_a^{\mathrm{lat}}(0)
  \sim
  \sum_{x\in\Lambda} x_a\,\mu_x,
  \qquad
  \mu_x = m\,n_x
  \quad\text{when the particles have common mass }m.
\]
This is a proposal-level lattice formula until the model supplies the density
\(\mu_x\), the conservation law, and the continuity/current relation.

\subsection{Compiler Implication}

The compiler rule is therefore stricter than in CA-16:
\[
\begin{array}{ll}
\text{Lorentzian candidate layer:}
  &(\Lambda,x,\{h_x\})\text{ may generate }K_a\sim\sum_x x_a h_x,\\[2mm]
\text{Galilean candidate layer:}
  &(\Lambda,x,\{h_x\})\text{ is not enough.}
\end{array}
\]
For Galilean boosts the input must also name:
\[
  \{\mu_x\}\text{ or }\{n_x\},\qquad
  M=\sum_x\mu_x,\qquad
  [H,M]=0,
\]
plus, for a momentum candidate, a current witness for the discrete continuity
equation.  If those data are absent, the compiler must fail loudly rather than
silently using the energy first moment.

\subsection{Why Projective Data Matter Earlier}

In the Hilbert-space grammar, projective symmetries were handled by passing to a
central extension.  Galilean symmetry is the first lattice target where that
choice becomes operational: the mass parameter is expected to sit in the
projective/central part of the quantum representation.  Locally we can support
only this precise statement:
\[
  \text{projective quantum symmetries are represented by unitary representations
  of central extensions.}
\]
The stronger statement identifying the Galilei extension by its standard
literature name remains a source-acquisition task.

\subsection{Boundary for CA-18--CA-22}

The follow-up block does four things:
\[
\begin{array}{ll}
\text{CA-19} & \text{derives the unextended Galilei vector-field algebra},\\
\text{CA-20} & \text{derives the mass central coefficient on a canonical core},\\
\text{CA-21} & \text{derives the lattice first-moment continuity formula},\\
\text{CA-22} & \text{turns the data into compiler rules and small examples}.
\end{array}
\]
No shard in this block claims that a finite lattice has exact Galilean symmetry
unless the corresponding finite-dimensional representation and commutator
checks are actually supplied.
